\chapter{Semantic Zoom}

\chapterepigraph{%
  Every level of description is correct.\\
  The question is not which level is true\\
  but which level is useful for the current question.
}{Flyxion, \textit{Semantic Infrastructure}}

\sidenote{Connects to\\renormalization\\tower (Ch.\,8)\\and constraint\\basins (Ch.\,24)}

Chapter~32 argued for a terrain model of knowledge. This chapter develops
the key feature of that terrain: its \emph{multi-scale structure}. The
same knowledge can be viewed at different resolutions — from the
fine-grained (individual propositions) to the coarse-grained (field-level
summaries). Semantic zoom is the operation that moves between these
resolutions.

\section{Scale-Dependent Representations}

Every domain of knowledge has a natural hierarchy of scales:

\begin{itemize}
  \item \textbf{Biology:} molecules $\to$ cells $\to$ tissues $\to$
        organisms $\to$ populations $\to$ ecosystems.
  \item \textbf{Physics:} quantum fields $\to$ particles $\to$ atoms
        $\to$ materials $\to$ systems $\to$ cosmological structures.
  \item \textbf{History:} events $\to$ episodes $\to$ periods $\to$
        eras $\to$ civilizational arcs.
  \item \textbf{Mathematics:} lemmas $\to$ theorems $\to$ theories
        $\to$ fields $\to$ mathematical traditions.
\end{itemize}

Each scale has its own vocabulary, its own causal mechanisms, and its
own productive questions. Confusing scales — asking cell-biology
questions at the organism level, or organism-level questions at the
ecosystem level — produces category errors.

\begin{definition}{Semantic Zoom Operator}{sem-zoom}
  A \emph{semantic zoom operator} is a family of projections
  $\{\pi_s : K \to K_s\}_{s \geq 0}$ where $K_s$ is the representation
  of knowledge $K$ at scale $s$. Zooming out ($s \to \infty$) is
  renormalization: detail is integrated out, coarse features emerge.
  Zooming in ($s \to 0$) is instantiation: abstract concepts are
  given specific content.
\end{definition}

\section{Micro, Meso, and Macro Views}

Three canonical scales appear repeatedly across domains:

\begin{itemize}
  \item \textbf{Micro:} the individual unit — atom, cell, event, lemma.
        The level of direct observation and experiment.
  \item \textbf{Meso:} the intermediate aggregate — molecule, tissue,
        episode, theorem. The level where most scientific activity occurs.
  \item \textbf{Macro:} the large-scale pattern — material, ecosystem,
        era, field. The level where global constraints and tendencies
        are visible.
\end{itemize}

Productive research typically moves between scales: a micro observation
(an anomalous data point) motivates a meso hypothesis (a new mechanism),
which predicts a macro pattern (a new large-scale regularity).

\section{Context-Sensitive Abstraction}

The appropriate scale for a representation depends on the question
being asked. This is \emph{context-sensitive abstraction}: the same
knowledge is represented at different scales depending on context.

\begin{definition}{Context-Optimal Scale}{ctx-optimal-scale}
  The \emph{context-optimal scale} for knowledge $K$ in context $C$
  is:
  \[
    s^*(K, C) = \arg\max_s I(K_s; \mathcal{Q}_C)
  \]
  where $I(K_s; \mathcal{Q}_C)$ is the mutual information between
  the scale-$s$ representation of $K$ and the set of questions
  $\mathcal{Q}_C$ relevant in context $C$.
\end{definition}

Experts have well-calibrated $s^*$: they know which scale of description
is appropriate for each question. Novices often operate at the wrong
scale — too much detail when a summary is needed, too much abstraction
when a specific example is needed.

\begin{namedthm}{The Multi-Scale Principle}
  Any knowledge system that must answer questions at multiple scales
  requires explicit semantic zoom capabilities: mechanisms for moving
  between scales efficiently, and for identifying the context-optimal
  scale for each query. Systems optimized for a single scale sacrifice
  performance at all other scales.
\end{namedthm}

\begin{exercise}
  Take a topic you know well and write three descriptions of it at
  three different scales (micro, meso, macro), each approximately
  100 words. Identify what information is present at each scale and
  absent at the others. What question is each scale best suited to answer?
\end{exercise}

\begin{exercise}
  The internet is a micro-scale knowledge system: individual pages
  are the primary unit, and large-scale structure must be inferred
  (through search, navigation, link-following). Design a ``macro-layer''
  for the internet that makes large-scale knowledge structure explicit.
  What would its data model look like?
\end{exercise}
